% !Mode:: "TeX:UTF-8"

\chapter{SoC 系统设计思路}

\section{设计目标与总体原则}

本课题 SoC 设计的核心目标是为乱序处理器提供可运行操作系统的完整系统环境，并在保证可调试性的前提下兼顾性能与实现复杂度。为此，系统采用“分层总线 + 多层次存储 + 可视化外设联调”的总体思路：高速访存路径优先保障带宽与时延，低速外设路径优先保障协议兼容与工程可维护性，同时通过统一验证框架降低系统集成阶段的调试成本。

在实现策略上，SoC 结构遵循模块化和可替换原则。处理器核、总线互连、存储控制器和外设控制器通过标准接口连接，使得各模块可以独立迭代。该设计既满足当前单发射乱序处理器的系统集成需求，也为后续微架构升级和外设扩展预留了接口空间。

\section{分层总线架构设计}

SoC 采用高级可扩展接口 4（AXI4, Advanced eXtensible Interface 4）\cite{math2012axi4} 与高级外设总线（APB, Advanced Peripheral Bus）组合的分层互连架构，以匹配不同设备的速率特征。处理器核心通过 AXI4 主接口接入一级 Crossbar，一级 Crossbar 重点服务主存访问关键路径，仅挂接高带宽存储控制器；二级 Crossbar 连接片上存储与桥接模块，承担中速设备访问；低速外设统一通过 AXI4ToAPB 桥进入 APB 域，由 APB Fanout 完成地址分发。

SoC 整体互连架构如图~\ref{fig:soc-arch} 所示。

\begin{figure}[H]
  \centering
  \includegraphics[width=0.95\textwidth]{soc-arch.png}
  \caption{SoC 系统整体互连架构}
  \label{fig:soc-arch}
\end{figure}

如图~\ref{fig:soc-arch} 所示，SoC 整体互连分为三个速率域。处理器核心作为唯一的 AXI4 Master 接入一级 AXI4 Crossbar。一级 Crossbar 的左侧分支经 AXI4 Delayer（用于在仿真中注入可配置的访存延迟）连接 SDRAM Controller，右侧分支经 AXI4 Fragmenter（将突发传输拆分为单拍事务以适配不支持 Burst 的从设备）与 AXI4 UserYanker（剥离用户自定义信号以对齐接口宽度）连接二级 AXI4 Crossbar。SDRAM 直接挂载在一级 Crossbar 上，以最短路径获得最高带宽。

二级 Crossbar 下挂 MROM（存放 FSBL, first stage boot loader）、SRAM（低延迟片上存储）以及 AXI4ToAPB 总线桥。总线桥将 AXI4 事务转换为 APB 协议后，由 APB Fanout 按地址空间分发至各低速外设：XIP Flash Controller、PSRAM QSPI Master、UART 16550、GPIO、PS2 键盘控制器、VGA 显示控制器，以及 PLIC（外部中断控制器）和 CLINT（定时器与软件中断控制器）。

在 Flash 访问路径上，XIP Flash Controller 与 APB Fanout 共享底层 SPI Controller 的物理访问权，通过 SPI Arbiter 进行仲裁：XIP 以高优先级发起 SPI 事务用于直接取指执行，APB 以低优先级提供 SPI 寄存器的直接读写通道。各外设控制器的 TX/RX 或数据端口连接至 NVBoard 虚拟开发板对应的外设仿真（串口终端、LED/数码管、键盘、VGA 显示器），以虚线表示。

该分层设计的关键收益在于：其一，高速域与低速域解耦，SDRAM 路径不受慢设备干扰；其二，AXI4 突发与并发事务能力被主存路径充分利用，APB 侧保持实现简单、硬件开销低；其三，桥接与仲裁逻辑集中管理，扩展新外设只需在 APB Fanout 中追加地址映射。

\section{多层次存储与启动链路设计}

存储系统采用 MROM、SRAM、Flash、PSRAM、SDRAM 的多层次组织。MROM 用于放置最小启动代码，SRAM 提供低延迟片上存储，Flash 提供非易失程序镜像并支持 XIP，PSRAM 与 SDRAM 提供更大容量主存。不同存储介质按“启动可靠性—访问延迟—容量”进行功能分工，形成从上电到系统运行的连续存储支撑。

在控制器实现上，关键 IP 核采用 Chisel~\cite{bachrach2012chisel} 实现或重写：SDRAM 控制器面向 AXI4 高带宽路径进行流水化组织，Flash/PSRAM 相关控制器面向 SPI/QSPI/QPI 协议兼容设计。Flash 访问同时支持寄存器方式与 XIP 方式，并通过仲裁机制共享底层 SPI 资源。该设计兼顾了启动阶段可执行性与运行阶段访存效率。

启动链路采用分级引导思路：上电后从 Flash 中的引导入口执行，完成基础初始化后将后续引导或系统镜像加载至 SDRAM，再跳转进入运行态。该流程减少了直接从串行 Flash 执行大体量程序的性能损失，为操作系统（如 xv6~\cite{cox2024xv6}、RT-Thread~\cite{yi2023rtthread}）启动提供更稳定的执行环境。

\begin{figure}[H]
  \centering
  \includegraphics{boot-flow.png}
  \caption{启动链路流程图}
  \label{fig:boot-flow}
\end{figure}

\section{外设集成与可视化验证思路}

外设侧围绕“可用、可测、可观测”进行设计。UART、GPIO、PS2、VGA 等控制器通过 APB 域接入，保证接口一致性与调试便利性。针对教学与开发场景，系统进一步集成 NVBoard，将串口回显、按键输入、数码管显示与 VGA 输出纳入同一仿真闭环，使 SoC 级功能验证从单纯日志比对扩展为可交互的行为验证。

在验证方法上，SoC 联调与指令级差分测试协同使用：差分测试用于保证体系结构状态正确性，NVBoard 与外设交互用于暴露系统级时序和接口问题，二者互补提升了复杂问题定位效率。基于该思路，系统能够稳定支撑操作系统运行与外设协同工作。

\section{本章小结}

本章从目标约束、总线分层、存储组织、启动链路与外设验证五个方面阐述了 SoC 设计思路。通过 AXI4/APB 分层互连、多层次存储分工、关键控制器 Chisel 化实现以及 NVBoard 可视化联调，系统在性能、可维护性与可验证性之间取得了平衡，为乱序处理器运行 xv6 等操作系统提供了完整平台基础。
